\origpage{353}{N. Robertson, P.D. Seymour / Journal of Combinatorial Theory, Series B 90 (2004) 325--385}
\hspace*{12pt}$\bullet$ \textit{for every countable sequence $(\Gamma_i, \phi_i, \mathscr{L}_i)$ $(i = 1, 2, \dots)$ of members of $\mathscr{P}$ there exist $j > i \ge 1$ and an inflation $\sigma$ of $M(\Gamma_i, \mathscr{L}_i)$ in $\Gamma_j$, such that}

\hspace*{24pt}$\circ$ \textit{for each $v \in V(M(\Gamma_i, \mathscr{L}_i)), V(\sigma(v)) \subseteq V(M(\Gamma_j, \mathscr{L}_j))$,}

\hspace*{24pt}$\circ$ \textit{for each $e \in E(M(\Gamma_i, \mathscr{L}_i)), \sigma(e) \in E(M(\Gamma_j, \mathscr{L}_j))$ and $\phi_i(e) \le \phi_j(\sigma(e))$,}

\hspace*{24pt}$\circ$ $|\bar{\pi}(\mathscr{L}_i)| = |\bar{\pi}(\mathscr{L}_j)|$, \textit{and for} $1 \le h \le |\bar{\pi}(\mathscr{L}_i)|$ \textit{if} $v$ \textit{is the} $h$th \textit{term of} $\pi(\mathscr{L}_i)$ \textit{then} $\sigma(v)$ \textit{contains the} $h$th \textit{term of} $\pi(\mathscr{L}_j)$.

\textit{Then $\mathscr{P}$ is well-behaved.}

The proof is easy\del{ and we omit it} (see [10] for the proofs of several similar results). \fix{We give the proof. Let $\Omega$ be a well-quasi-order, let $(\Gamma_i, \phi_i, \mathscr{L}_i)$ $(i = 1, 2, \dots)$ be members of $\mathscr{P}$, and let $\xi_i \colon \mathscr{L}_i \to E(\Omega)$ for each $i$. For each $i$ let $A^i_1, \dots, A^i_{k_i}$ be the members of $\mathscr{L}_i$ in some order, where $k_i \le m$; and for $1 \le t \le k_i$ let $c_i(t)$ be the sequence of the positions in $\pi(\mathscr{L}_i)$ of the terms of $\pi(A^i_t)$, in order. (Each term of $\pi(A^i_t)$ is a term of $\pi(\mathscr{L}_i)$, since $\bar{\pi}(\mathscr{L}_i) = \bigcup(\bar{\pi}(A) : A \in \mathscr{L}_i)$.) Since $k_i \le m$, and each $c_i(t)$ is a sequence of at most $n$ numbers in $\{1, \dots, n\}$, the sequence $(c_i(1), \dots, c_i(k_i))$ takes only finitely many values; so by passing to an infinite subsequence, we may assume that $k_i = k$ and $c_i(t) = c(t)$ for all $i \ge 1$ and $1 \le t \le k$. Since every sequence of members of $E(\Omega)$ has an infinite subsequence in which each term is at most every later one, by passing to a subsequence $k$ times we may also assume that $\xi_i(A^i_t) \le \xi_j(A^j_t)$ for all $j > i \ge 1$ and $1 \le t \le k$. By the second hypothesis there exist $j > i \ge 1$ and an inflation $\sigma$ of $M(\Gamma_i, \mathscr{L}_i)$ in $\Gamma_j$ satisfying its three conditions. Let $\tau(A^i_t) = A^j_t$ for $1 \le t \le k$. Then $\xi_i(A) \le \xi_j(\tau(A))$ for all $A \in \mathscr{L}_i$, and so it suffices to show that $\tau$ is an outline of $(\Gamma_i, \phi_i, \mathscr{L}_i)$ in $(\Gamma_j, \phi_j, \mathscr{L}_j)$.

The first and last conditions in the definition of an outline hold by the second condition on $\sigma$. The map $\tau$ is injective since $A^j_1, \dots, A^j_k$ are distinct, and $|\bar{\pi}(\tau(A^i_t))| = |\bar{\pi}(A^i_t)|$ since both equal the length of $c(t)$. For the third condition, let $A = A^i_t$, and let $v$ be the $h$th term of $\pi(A)$. Then $v$ is the $c$th term of $\pi(\mathscr{L}_i)$, where $c$ is the $h$th term of $c(t)$, and so $\sigma(v)$ contains the $c$th term of $\pi(\mathscr{L}_j)$ by the third condition on $\sigma$; and since $c_j(t) = c(t)$, this is the $h$th term of $\pi(\tau(A))$. For the fourth condition, let $A \in \mathscr{L}_i$ and $v \in V(M(\Gamma_i, \mathscr{L}_i))$, and let $x, y \in V(\sigma(v)) \cap V(\tau(A))$. By the first condition on $\sigma$, $x, y \in V(M(\Gamma_j, \mathscr{L}_j))$, and so $x, y \in \bar{\pi}(\tau(A))$, since $M(\Gamma_j, \mathscr{L}_j) \subseteq \tau(A)^c$ and $V(\tau(A)^- \cap \tau(A)^c) = \bar{\pi}(\tau(A))$. Let $x, y$ be the $h$th and $h'$th terms of $\pi(\tau(A))$, and let $u, u'$ be the $h$th and $h'$th terms of $\pi(A)$. By what we just proved, $x \in V(\sigma(u))$ and $y \in V(\sigma(u'))$. Since the subgraphs $\sigma(w)$ $(w \in V(M(\Gamma_i, \mathscr{L}_i)))$ are mutually disjoint, it follows that $u = v = u'$; so $h = h'$, and $x = y$. Hence $|V(\sigma(v)) \cap V(\tau(A))| \le 1$, and so $\tau$ is an outline. This proves 9.3.}\fixnote{fix:9.3-proof-added}{The paper omits the proof of 9.3. The argument above is the one used in the formalization.} We define the image of a $\chi$-place under a $\Phi_\chi$-preserving homeomorphism in the natural way.

\textbf{9.4.} \textit{Let $\chi$ be a disjoint colour scheme, and let $\mathscr{C}$ be a well-behaved set of $\chi$-places. Let $\mathscr{C}'$ be the set of all images of members of $\mathscr{C}$ under $\Phi_\chi$-preserving homeomorphisms of $\Sigma_\chi$. Then $\mathscr{C}'$ is well-behaved.}

Again, the proof is clear.

\sectionhead{10. Tangle flaws}

Our next objective is to produce some well-behaved sets of $\chi$-places. The proofs of several of these theorems are similar, and so we give only the first in detail. Throughout this section and the next, $\chi$ is a disjoint colour scheme, and $\mathscr{S}$ is a similarly oriented set of $\chi$-coloured paintings, such that if $(\Gamma, \phi) \in \mathscr{S}$ and $\alpha$ is a $\Phi_\chi$-preserving homeomorphism of $\Sigma_\chi$, then $\mathscr{S}$ contains the image of $(\Gamma, \phi)$ under $\alpha$ (briefly, $\mathscr{S}$ is closed under $\Phi_\chi$-preserving homeomorphisms). Let $\mathscr{D}$ be the set of all quadruples $(\Gamma, \phi, \lambda, \mathscr{T})$ such that $(\Gamma, \phi) \in \mathscr{S}$, $\lambda$ is a tie-breaker in $\Gamma$, and $\mathscr{T}$ is a tangle in $\Gamma$.

Let $(\Gamma, \phi, \lambda, \mathscr{T}) \in \mathscr{D}$, such that $rep(\Gamma), ord(\mathscr{T}) > rep(\mathscr{T})$. Then $\hat{\Sigma}_\chi$ is not a sphere, and there is a circuit $F$ of $\Gamma^*$ with $\frac{1}{2}|E(F)| \le rep(\mathscr{T})$ such that $\Sigma_1$ is not a disc, where $\Sigma_1, \Sigma_2$ are the closures of the two components of $\hat{\Sigma}_\chi \setminus U(F)$ and $(\Gamma \cap \Sigma_1, \Gamma \cap \Sigma_2) \in \mathscr{T}$. (It follows that $\frac{1}{2}|E(F)| = rep(\mathscr{T})$ for every such $F$.) Let us choose such a circuit $F$ with $\lambda(F)$ minimal. We call $F$ a \textit{representativeness flaw} for $(\Gamma, \phi, \lambda, \mathscr{T})$.

\textbf{10.1.} \textit{With $\Gamma, \phi, \lambda, \mathscr{T}, F, \Sigma_1, \Sigma_2$ as above,}

\hspace*{12pt}$\bullet$ \textit{$\Sigma_2$ is a disc,}

\hspace*{12pt}$\bullet$ \textit{either $|U(F) \cap bd(\Sigma_\chi)| \le 1$ and $U(F) \subseteq \Sigma_\chi$, or $|U(F) \cap bd(\Sigma_\chi)| = 2$ and $U(F) \not\subseteq \Sigma_\chi$, or $dist(\Gamma) \le \frac{1}{4}|E(F)| + 1$,}

\hspace*{12pt}$\bullet$ \textit{$F$ is $\lambda$-linked to $\mathscr{T}$.}
